\subsection{Inmersiones regulares y cuasirregulares.}
\label{IV.16.9}

\begin{definition}[16.9.1]
\label{IV.16.9.1}
Sea $X$ un espacio anillado. Se dice que un ideal $\mathscr{I}$ de
$\mathcal{O}_X$ es \emph{regular} (resp. \emph{cuasirregular}) si, para todo
punto $x\in\operatorname{Supp}(\mathcal{O}_X/\mathscr{I})$, existen una
vecindad abierta $U$ de $x$ en $X$ y una sucesión regular (0, 15.2.2) (resp.
cuasirregular (0, 15.2.2)) de elementos de $\Gamma(U,\mathcal{O}_X)$ que
engendra $\mathscr{I}|U$.
\end{definition}

Se dirá que una sucesión regular (resp. cuasirregular) de secciones de
$\mathcal{O}_X$ sobre $U$ que engendra $\mathscr{I}|U$ es un
\emph{sistema regular} (resp. \emph{cuasirregular}) \emph{de generadores}
de $\mathscr{I}|U$.

\begin{definition}[16.9.2]
\label{IV.16.9.2}
Sea $j:Y\to X$ una inmersión de preesquemas y sea $U$ un abierto de $X$ tal
que $j(Y)\subset U$ y que $j$ sea una inmersión cerrada de $Y$ en $U$. Se
dice que $j$ es \emph{regular} (resp. \emph{cuasirregular}) si el
subpreesquema cerrado $j(Y)$ de $U$ asociado a $j$ está definido por un ideal
regular (resp. cuasirregular) de $\mathcal{O}_U$ (condición independiente del
abierto $U$ elegido).
\end{definition}

Se dice que un subpreesquema $Y$ de un preesquema $X$ está \emph{inmerso
regularmente} (resp. \emph{inmerso cuasirregularmente}) si la inyección
canónica $j:Y\to X$ es una inmersión regular (resp. cuasirregular). Si $Y$ es
un subpreesquema cerrado de $X$ y $\mathscr{I}$ el ideal de $\mathcal{O}_X$
que define $Y$, esto equivale a decir que $\mathscr{I}$ es \emph{regular}
(resp. \emph{cuasirregular}).

Por ejemplo, si $A$ es un anillo \emph{íntegro} y $f$ un elemento $\neq0$ de
$A$, el subpreesquema cerrado $V(f)$ de $\operatorname{Spec}(A)$ (isomorfo a
$\operatorname{Spec}(A/fA)$) está \emph{inmerso regularmente} en
$\operatorname{Spec}(A)$.

Todo ideal regular es cuasirregular (0, 15.2.2); toda inmersión regular es
cuasirregular (cf. (16.9.11) para un enunciado recíproco).

\begin{proposition}[16.9.3]
\label{IV.16.9.3}
Sean $X$ un espacio anillado, $\mathscr{I}$ un ideal de $\mathcal{O}_X$ y
$(f_i)_{1\leq i\leq m}$ una sucesión finita de secciones de $\mathcal{O}_X$
sobre $X$ que engendra $\mathscr{I}$. Para que $(f_i)$ sea una sucesión
cuasirregular (0, 15.2.2), es necesario y suficiente que se satisfagan las
condiciones siguientes:
\begin{enumerate}
  \item[(i)] Las imágenes canónicas de los $f_i$ en
  $\mathscr{I}/\mathscr{I}^2$ forman una base de este
  $\mathcal{O}_X/\mathscr{I}$-módulo.
  \item[(ii)] El homomorfismo sobreyectivo canónico (16.1.2.2)
  \[
    \mathbf{S}_{\mathcal{O}_X/\mathscr{I}}^{\bullet}
    (\mathscr{I}/\mathscr{I}^2)
    \longrightarrow\mathscr{G}r_{\mathscr{I}}^{\bullet}(\mathcal{O}_X)
  \]
  es biyectivo.
\end{enumerate}
Además, si así ocurre, toda sucesión $(f'_i)_{1\leq i\leq m}$ de $m$
secciones de $\mathscr{I}$ sobre $X$ que engendra $\mathscr{I}$ es
cuasirregular.\footnote{El texto impreso usa $n$ en ambos lugares; debe ser
$m$, el número de elementos de la base dada por $(f_i)_{1\leq i\leq m}$. La
demostración exige que la nueva familia tenga el mismo número de elementos que
una base de $\mathscr{I}/\mathscr{I}^2$.}
\end{proposition}

\oldpage[IV-4]{47}

Las dos condiciones del enunciado no hacen sino expresar la definición dada
en (0, 15.2.2), teniendo en cuenta la definición de los homomorfismos
canónicos (0, 15.2.1.1). La última afirmación resulta de que, si un módulo $M$
sobre un anillo conmutativo $A$ admite una base de $n$ elementos, todo sistema
de generadores de $M$ que tenga $n$ elementos es una base de $M$ (Bourbaki,
\emph{Alg. comm.}, cap.~II, §~3, cor.~5 del teorema~1).

\begin{corollary}[16.9.4]
\label{IV.16.9.4}
Sean $X$ un espacio localmente anillado y $\mathscr{I}$ un ideal de
$\mathcal{O}_X$. Para que $\mathscr{I}$ sea cuasirregular, es necesario y
suficiente que satisfaga las condiciones siguientes:
\begin{enumerate}
  \item[(i)] $\mathscr{I}$ es de tipo finito.
  \item[(ii)] $\mathscr{I}/\mathscr{I}^2$ es un
  $(\mathcal{O}_X/\mathscr{I})$-módulo localmente libre.
  \item[(iii)] El homomorfismo canónico
  \[
    \mathbf{S}_{\mathcal{O}_X/\mathscr{I}}^{\bullet}
    (\mathscr{I}/\mathscr{I}^2)
    \longrightarrow\mathscr{G}r_{\mathscr{I}}^{\bullet}(\mathcal{O}_X)
    \tag{16.9.4.1}
    \label{IV.16.9.4.1}
  \]
  es biyectivo.
\end{enumerate}
\end{corollary}

La necesidad de las condiciones resulta inmediatamente de (16.9.3). Para ver
que son suficientes, hay que demostrar, en virtud de (16.9.3), que, si en un
punto $x\in\operatorname{Supp}(\mathcal{O}_X/\mathscr{I})$ existen una
vecindad abierta $U$ de $x$ en $X$ y $n$ secciones $f_i$
($1\leq i\leq n$) de $\mathscr{I}$ sobre $U$ cuyas imágenes canónicas en
$\mathscr{I}/\mathscr{I}^2$ forman una base de
$(\mathscr{I}/\mathscr{I}^2)|U$ sobre
$(\mathcal{O}_X/\mathscr{I})|U$, entonces existe una vecindad abierta
$V\subset U$ de $x$ tal que los $f_i|V$ engendran $\mathscr{I}|V$. Ahora
bien, por hipótesis, $\mathscr{I}_x\neq\mathcal{O}_x$, por lo que
$\mathscr{I}_x$ está contenido en el ideal maximal de $\mathcal{O}_x$; como
$\mathscr{I}_x$ es un $\mathcal{O}_x$-módulo de tipo finito y las clases de
los $(f_i)_x$ en $\mathscr{I}_x/\mathscr{I}_x^2$ engendran este
$(\mathcal{O}_x/\mathscr{I}_x)$-módulo, el lema de Nakayama muestra que los
$(f_i)_x$ engendran $\mathscr{I}_x$. Como $\mathscr{I}$ es de tipo finito,
se concluye por (0\textsubscript{I}, 5.2.2).

\begin{corollary}[16.9.5]
\label{IV.16.9.5}
Sean $X$ un espacio localmente anillado, $\mathscr{I}$ un ideal cuasirregular
de $\mathcal{O}_X$, $(f_i)_{1\leq i\leq n}$ una sucesión de secciones de
$\mathscr{I}$ sobre $X$, y
$x$ un punto de $\operatorname{Supp}(\mathcal{O}_X/\mathscr{I})$. Las
condiciones siguientes son equivalentes:
\begin{enumerate}
  \item[(a)] Existe una vecindad abierta $U$ de $x$ en $X$ tal que los
  $f_i|U$ forman una sucesión cuasirregular de elementos de
  $\Gamma(U,\mathcal{O}_X)$ que engendra $\mathscr{I}|U$.
  \item[(b)] Los $(f_i)_x$ forman un sistema de generadores de
  $\mathscr{I}_x$ cuyo número de elementos es el menor posible.
  \item[(b')] Los $(f_i)_x$ forman un sistema minimal de generadores de
  $\mathscr{I}_x$.
  \item[(c)] Si $\bar f_i$ es la imagen canónica de $f_i$ en
  $\Gamma(X,\mathscr{I}/\mathscr{I}^2)$, los $(\bar f_i)_x$ forman una base
  del $(\mathcal{O}_x/\mathscr{I}_x)$-módulo
  $\mathscr{I}_x/\mathscr{I}_x^2$.
\end{enumerate}
\end{corollary}

Por hipótesis, $\mathcal{O}_x$ es un anillo local y $\mathscr{I}_x$ un ideal
de tipo finito de $\mathcal{O}_x$ contenido en el ideal maximal de
$\mathcal{O}_x$; la equivalencia entre b), b') y c) resulta, por tanto, del
lema de Nakayama (Bourbaki, \emph{Alg. comm.}, cap.~II, §~3,
n\textsuperscript{o}~2, prop.~5). Es claro que a) implica c) en virtud de
(16.9.3); por otra parte, de (0\textsubscript{I}, 5.2.2) resulta que, si se
satisface la condición c) (y, por tanto, también b)), existe una vecindad $U$
de $x$ en $X$ tal que $(\mathscr{I}/\mathscr{I}^2)|U$ tiene rango constante
igual a $n$, y que los $f_i|U$ engendran $\mathscr{I}|U$; basta, por tanto,
aplicar en $U$ la última afirmación de (16.9.3).

\begin{remarks}[16.9.6]
\label{IV.16.9.6}
\begin{enumerate}
  \item[(i)] Bajo las hipótesis generales de (16.9.5), para que la sucesión
  $(f_i)$ engendre $\mathscr{I}$ no basta con que los $(\bar f_i)_y$ formen
  una base del $(\mathcal{O}_y/\mathscr{I}_y)$-módulo
  $\mathscr{I}_y/\mathscr{I}_y^2$ para todo $y\in X$. Se obtiene un ejemplo
  tomando
\oldpage[IV-4]{48}
  $X=\operatorname{Spec}(A)$, donde $A$ es un anillo de Dedekind, y
  $\mathscr{I}=\widetilde{\mathfrak{I}}$, donde $\mathfrak{I}$ es un ideal
  primo no principal de $A$; se tiene entonces, en efecto,
  $\mathscr{I}_y/\mathscr{I}_y^2=0$ en todo punto $y$ distinto del punto
  $x\in X$ correspondiente a $\mathfrak{I}$, y
  $\mathscr{I}_x/\mathscr{I}_x^2$ es de rango $1$ sobre el cuerpo
  $\mathcal{O}_x/\mathscr{I}_x$; además, $\mathscr{I}$ es evidentemente un
  ideal regular.
  \item[(ii)] En (16.9.5) no se puede sustituir «cuasirregular» por
  «regular», ni siquiera cuando $X$ es un preesquema (cf. (16.9.12)). En
  efecto, designemos por $B$ el anillo de los gérmenes de funciones
  indefinidamente diferenciables en el punto $0$ de $\mathbb{R}$; tiene un
  ideal maximal $\mathfrak{m}$ engendrado por el germen $t$ de la aplicación
  identidad de $\mathbb{R}$ en el punto $0$, y la intersección $\mathfrak{n}$
  de los $\mathfrak{m}^k$ para $k>0$ no se reduce a $0$. Sea ahora $A$ el
  anillo cociente $B[T]/\mathfrak{n}TB[T]$, y sean $f_1,f_2$ las imágenes
  canónicas en $A$ de los elementos $t$ y $T$ de $B[T]$. La sucesión
  $(f_1,f_2)$ es regular en $A$: en efecto, $f_1$ no es divisor de $0$ en
  $A$, pues la relación $tP[T]\in\mathfrak{n}TB[T]$, para un polinomio
  $P\in B[T]$, implica que los productos de $t$ por los coeficientes de $P$
  pertenecen al ideal $\mathfrak{n}$, de lo que resulta inmediatamente que
  estos coeficientes están a su vez en $\mathfrak{n}$, y, por tanto,
  $P[T]\in\mathfrak{n}TB[T]$. Como $B/tB$ es isomorfo a $\mathbb{R}$,
  $A/f_1A$ es isomorfo al anillo de polinomios $\mathbb{R}[T]$, y, por tanto,
  es íntegro; la imagen de $f_2$ en $A/f_1A$, al ser igual a $T$, no es, por
  consiguiente, divisor de $0$, lo que demuestra la afirmación. Sin embargo,
  $f_2$ es divisor de $0$ en $A$, pues para todo elemento no nulo
  $x\in\mathfrak{n}$, la imagen de $x$ en $A$ es $\ne0$, mientras que la
  imagen de $xT$ es nula. Se concluye que la sucesión $(f_2,f_1)$ no es
  regular en $A$; por otra parte, el ideal $\mathfrak{I}=f_1A+f_2A$ es
  distinto de $A$, por lo que las condiciones b), b') y c) de (16.9.5) no
  implican la condición a) cuando en ella se sustituye «cuasirregular» por
  «regular».
\end{enumerate}
\end{remarks}

\begin{env}[16.9.7]
\label{IV.16.9.7}
Si $X=\operatorname{Spec}(A)$ es un esquema afín, se dirá que un ideal
$\mathfrak{I}$ de $A$ es regular (resp. cuasirregular) si el ideal
$\mathscr{I}=\widetilde{\mathfrak{I}}$ de $\mathcal{O}_X$ es regular
(resp. cuasirregular); obsérvese que esta noción es local y no implica en
absoluto la existencia de un sistema de generadores de $\mathfrak{I}$ que
forme en $A$ una sucesión regular (resp. cuasirregular), como muestra el
ejemplo (16.9.6, (i)); sin embargo, así ocurre cuando $A$ es local (16.9.5).

La proposición (16.9.4) se expresa en términos de inmersiones cuasirregulares
de la manera siguiente:
\end{env}

\begin{proposition}[16.9.8]
\label{IV.16.9.8}
Sea $j:Y\to X$ un morfismo de preesquemas; para que $j$ sea una inmersión
cuasirregular, es necesario y suficiente que $j$ satisfaga las condiciones
siguientes:
\begin{enumerate}
  \item[(i)] $j$ es una inmersión localmente de presentación finita.
  \item[(ii)] El haz conormal
  $\operatorname{Gr}^1(j)=\mathcal{N}_{Y/X}$ (16.1.2) es un
  $\mathcal{O}_Y$-módulo localmente libre.
  \item[(iii)] El homomorfismo canónico
  \[
    \operatorname{S}_{\mathcal{O}_Y}^{\bullet}
      (\operatorname{Gr}^1(j))\longrightarrow\operatorname{Gr}^{\bullet}(j)
  \]
  (16.1.2.2) es biyectivo.
\end{enumerate}
\end{proposition}

Como la cuestión es local sobre $Y$, se puede reducir al caso en que $j$ es
la inyección canónica de un subpreesquema cerrado $Y$ de $X$, y entonces la
expresión de (16.9.4) en la forma (16.9.8) resulta de la explicitación de
$\operatorname{Gr}^1(j)$ y $\operatorname{Gr}^{\bullet}(j)$ en términos del
ideal $\mathscr{I}$ de $\mathcal{O}_X$ que define el subpreesquema $Y$
(16.1.3, (ii)).

\begin{corollary}[16.9.9]
\label{IV.16.9.9}
Sean $Y$ un preesquema, $X$ un $Y$-preesquema y $j:Y\to X$ una $Y$-sección
de $X$, de modo que el $n$-ésimo invariante normal $\mathscr{A}^{(n)}$ de $j$
(16.1.2) es una $\mathcal{O}_Y$-álgebra aumentada (16.1.7); pongamos
$\mathscr{A}^{(\infty)}=\varprojlim\mathscr{A}^{(n)}$. Para que $j$ sea una
inmersión cuasirregular, es necesario y suficiente que $j$ sea localmente de
presentación finita y que todo $y\in Y$ admita una vecindad abierta afín $U$
cuyo anillo es $C$ tal que $\mathscr{A}^{(\infty)}|U$ sea isomorfa, como
$\mathcal{O}_U$-álgebra topológica aumentada, a
$\mathcal{O}_U[[T_1,\ldots,T_n]]$.
\end{corollary}

Se puede reducir al caso en que $j$ es una inmersión cerrada restringiéndose
a una vecindad suficientemente pequeña de $y$ (véase el razonamiento de
(16.4.11)), y entonces $\mathcal{O}_Y$ se identifica con el álgebra cociente
$\mathcal{O}_X/\mathscr{I}$ y el homomorfismo sobreyectivo canónico
$\mathcal{O}_X\to\mathcal{O}_Y$ admite un
\oldpage[IV-4]{49}
inverso por la derecha (16.1.7). Se puede, por tanto, suponer que
$X=\operatorname{Spec}(B)$ e $Y=\operatorname{Spec}(A)$ son afines, que $B$
es una $A$-álgebra aumentada y que el ideal de aumento $\mathfrak{I}$ es de
tipo finito. Como $\mathscr{A}^{(n)}$ se identifica entonces con
$(B/\mathfrak{I}^{n+1})^{\sim}$, el corolario resulta de la equivalencia entre
b) y c) en (0, 19.5.4), puesto que $B/\mathfrak{I}=A$.

Obsérvese que, en el caso afín considerado, el hecho de que $j$ sea una
inmersión cuasirregular equivale también, en virtud de (0, 19.5.4), a decir
que $B$ es una $A$-álgebra formalmente lisa para la topología
$\mathfrak{I}$-preádica.

Obsérvese asimismo que la condición de que $j$ sea una inmersión localmente de
presentación finita se satisface siempre cuando el morfismo $X\to Y$ es
localmente de tipo finito (1.4.3, (v)).

\begin{proposition}[16.9.10]
\label{IV.16.9.10}
Sean $X$ un preesquema localmente noetheriano, $Y$ un subpreesquema de $X$,
$j:Y\to X$ la inyección canónica y $y$ un punto de $Y$.
\begin{enumerate}
  \item[(i)] Para que exista una vecindad abierta $U$ de $y$ en $X$ tal que
  la restricción $Y\cap U\to U$ de $j$ sea una inmersión regular, es necesario
  y suficiente que el núcleo $\mathscr{I}_y$ del homomorfismo sobreyectivo
  $\mathcal{O}_{X,y}\to\mathcal{O}_{Y,y}$ esté engendrado por una sucesión
  regular de elementos de $\mathcal{O}_{X,y}$.
  \item[(ii)] Para que la inmersión $j$ sea regular, es necesario y suficiente
  que sea cuasirregular.
\end{enumerate}
\end{proposition}

\begin{enumerate}
  \item[(i)] Se puede reducir al caso en que $Y$ es un subpreesquema cerrado
  de $X$ definido por un ideal \emph{coherente} $\mathscr{I}$ de
  $\mathcal{O}_X$. La condición es evidentemente necesaria. Recíprocamente,
  si $\mathscr{I}_y$ está engendrado por una sucesión regular $(s_i)_y$, donde
  los $s_i$ son secciones de $\mathscr{I}$ sobre una vecindad abierta $U$ de
  $y$ en $X$, se puede suponer que los $s_i$ engendran $\mathscr{I}|U$
  (0\textsubscript{I}, 5.2.2) y forman una sucesión regular (0, 15.2.4), de
  donde se sigue la afirmación.
  \item[(ii)] El hecho de que una inmersión cuasirregular sea regular resulta
  de (i) y de la identidad entre las sucesiones cuasirregulares y las
  sucesiones regulares de $\mathcal{O}_{X,y}$ formadas por elementos del ideal
  maximal (0, 15.1.11).
\end{enumerate}

Cuando (sin hipótesis noetheriana sobre $X$) el núcleo $\mathscr{I}_y$ de
$\mathcal{O}_{X,y}\to\mathcal{O}_{Y,y}$ está engendrado por una sucesión
regular de elementos de $\mathcal{O}_{X,y}$, se dice que la inmersión $j$ es
\emph{regular en el punto $y$}.

\begin{corollary}[16.9.11]
\label{IV.16.9.11}
Sea $X$ un preesquema localmente noetheriano; entonces todo ideal
cuasirregular de $\mathcal{O}_X$ es regular.
\end{corollary}

\begin{remarks}[16.9.12]
\label{IV.16.9.12}
\begin{enumerate}
  \item[(i)] Obsérvese que una inmersión regular no es, en general, un
  morfismo plano, ni \emph{a fortiori} un morfismo \emph{regular} en el sentido
  de (6.8.1).
  \item[(ii)] Sea $A$ un anillo local noetheriano; de (16.9.4) y (0, 17.1.1)
  resulta inmediatamente que, para que $A$ sea \emph{regular}, es necesario y
  suficiente que su ideal maximal $\mathfrak{m}$ sea \emph{cuasirregular}
  (o \emph{regular}, lo que equivale a lo mismo puesto que $A$ es
  noetheriano). Para que un esquema afín noetheriano $X$ sea \emph{regular},
  es necesario y suficiente que, para todo punto cerrado $x\in X$, la
  inyección canónica $\operatorname{Spec}(k(x))\to X$ sea una inmersión
  \emph{regular}.
\end{enumerate}
\end{remarks}

\begin{proposition}[16.9.13]
\label{IV.16.9.13}
Sean $X$ un preesquema localmente noetheriano, $Y$ un subpreesquema de $X$ e
$Y'$ un subpreesquema de $Y$, tal que la inyección canónica $j:Y'\to Y$ sea
regular. Entonces la sucesión de $\mathcal{O}_{Y'}$-módulos
\[
  \xymatrix{
    0\ar[r] & j^*(\mathcal{N}_{Y/X})\ar[r]
      & \mathcal{N}_{Y'/X}\ar[r] & \mathcal{N}_{Y'/Y}\ar[r] & 0
  }
  \tag{16.9.13.1}
  \label{IV.16.9.13.1}
\]
\oldpage[IV-4]{50}
es exacta; además, para todo $x\in X$, existe una vecindad abierta $U$ de $x$
tal que las restricciones a $U$ de los homomorfismos de (16.9.13.1) forman
una sucesión exacta y escindida.
\end{proposition}

Demostremos primero el lema siguiente:

\begin{lemma}[16.9.13.2]
\label{IV.16.9.13.2}
Sean $A$ un anillo, $\mathfrak{I}$ un ideal de $A$,
$A'=A/\mathfrak{I}$, $(f_i)_{1\leq i\leq r}$ una sucesión de elementos de
$A$ que es $A'$-regular, $\mathfrak{K}=\sum_i f_iA$,
$\mathfrak{L}=\mathfrak{I}+\mathfrak{K}$ y
$\mathfrak{K}'=\sum_i f_iA'$, de modo que $C=A/\mathfrak{L}$ es isomorfo a
$A'/\mathfrak{K}'$. Entonces, para todo entero $n>0$ y todo entero $N\geq n$,
se tiene la relación
\[
  \mathfrak{I}\cap\mathfrak{K}^n
  =\mathfrak{I}\mathfrak{K}^n+(\mathfrak{I}\cap\mathfrak{K}^N).
  \tag{16.9.13.3}
  \label{IV.16.9.13.3}
\]
\end{lemma}

Evidentemente, basta demostrar que todo elemento del primer miembro está
contenido en el segundo y, por inducción sobre $N-n$,\footnote{El texto
impreso atribuye la inducción a $n$; la iteración que reduce el enunciado al
caso $N=n+1$ es una inducción sobre la diferencia $N-n$.} se reduce al caso en
que $N=n+1$. Un elemento del primer miembro de (16.9.13.3), al pertenecer a
$\mathfrak{K}^n$, se escribe $P(f_1,\ldots,f_r)$, donde
$P\in A[T_1,\ldots,T_r]$ es homogéneo y de grado $n$. Si $f_i'$ es la imagen
canónica de $f_i$ en $A'$, la hipótesis
$P(f_1,\ldots,f_r)\in\mathfrak{I}$ significa que
$P(f_1',\ldots,f_r')=0$. Pero
$P(f_1',\ldots,f_r')\in\mathfrak{K}'^n$, por lo que la imagen canónica de
$P(f_1',\ldots,f_r')$ en
$\mathfrak{K}'^n/\mathfrak{K}'^{n+1}$ es nula. Ahora bien, la hipótesis de que
la sucesión $(f_i)$ es $A'$-regular implica que el homomorfismo canónico
\[
  \mathbf{S}_C^n(\mathfrak{K}'/\mathfrak{K}'^2)
  \longrightarrow\mathfrak{K}'^n/\mathfrak{K}'^{n+1}
\]
es biyectivo (0, 15.1.9); se concluye que los coeficientes de $P$ pertenecen
a $\mathfrak{L}=\mathfrak{I}+\mathfrak{K}$. De ello resulta inmediatamente
que $P(f_1,\ldots,f_r)\in\mathfrak{I}\mathfrak{K}^n+\mathfrak{K}^{n+1}$, y,
como $P(f_1,\ldots,f_r)\in\mathfrak{I}$, se tiene finalmente
\[
  P(f_1,\ldots,f_r)
  \in\mathfrak{I}\mathfrak{K}^n
    +(\mathfrak{I}\cap\mathfrak{K}^{n+1}),
\]
lo que demuestra el lema.

Tomando los cocientes de ambos miembros de (16.9.13.3) por
$\mathfrak{I}\mathfrak{K}^n$, se ve que las relaciones (16.9.13.3) para
$N\geq n$ implican
\[
  (\mathfrak{I}\cap\mathfrak{K}^n)/\mathfrak{I}\mathfrak{K}^n
  \subset\bigcap_{N\geq n}\mathfrak{K}^N
    \cdot(A/(\mathfrak{I}\mathfrak{K}^n)).
  \tag{16.9.13.4}
  \label{IV.16.9.13.4}
\]
De ello se deduce el

\begin{corollary}[16.9.13.5]
\label{IV.16.9.13.5}
Supongamos satisfechas las hipótesis de (16.9.13.2) y supongamos, además, que
el anillo $A$ es noetheriano y que $\mathfrak{K}$ está contenido en el radical
de $A$. Entonces, para todo entero $n>0$, se tiene
\[
  \mathfrak{I}\cap\mathfrak{K}^n=\mathfrak{I}\mathfrak{K}^n.
  \tag{16.9.13.6}
  \label{IV.16.9.13.6}
\]
\end{corollary}

En efecto, el segundo miembro de (16.9.13.4) es entonces nulo, puesto que
$A/\mathfrak{I}\mathfrak{K}^n$ es un $A$-módulo de tipo finito (Bourbaki,
\emph{Alg. comm.}, cap.~III, §~3, n\textsuperscript{o}~3, prop.~6).

Tomemos, en particular, $n=2$ en (16.9.13.6), y observemos que se tiene
$\mathfrak{L}^2=\mathfrak{I}^2+\mathfrak{I}\mathfrak{K}+\mathfrak{K}^2
=\mathfrak{I}\mathfrak{L}+\mathfrak{K}^2$; puesto que
$\mathfrak{I}\mathfrak{L}\subset\mathfrak{L}^2$, se deduce que
\[
  \mathfrak{I}\cap\mathfrak{L}^2
  =\mathfrak{I}\mathfrak{L}+(\mathfrak{I}\cap\mathfrak{K}^2)
  =\mathfrak{I}\mathfrak{L}+\mathfrak{I}\mathfrak{K}^2
  =\mathfrak{I}\mathfrak{L},
\]
es decir,
\[
  \mathfrak{I}\cap\mathfrak{L}^2=\mathfrak{I}\mathfrak{L},
  \tag{16.9.13.7}
  \label{IV.16.9.13.7}
\]
lo que también puede expresarse diciendo que el homomorfismo canónico
\[
  \mathfrak{I}/\mathfrak{I}\mathfrak{L}
  \longrightarrow(\mathfrak{I}+\mathfrak{L}^2)/\mathfrak{L}^2
\]
es biyectivo.

Demostrados estos lemas, demostremos la primera afirmación de (16.9.13): basta
\oldpage[IV-4]{51}
evidentemente demostrar que la sucesión de los tallos, en un punto $x\in Y'$,
de los haces que figuran en (16.9.13.1) es exacta. Ahora bien, si se pone
$A=\mathcal{O}_{X,x}$, se puede escribir
$\mathcal{O}_{Y,x}=A'=A/\mathfrak{I}$, donde $\mathfrak{I}$ es un ideal
contenido en el ideal maximal de $A$, y luego
$\mathcal{O}_{Y',x}=A'/\mathfrak{K}'$, donde $\mathfrak{K}'$ está engendrado
por una sucesión $A'$-regular de elementos de $A'$, que son a su vez imágenes
de los elementos de una sucesión $A'$-regular de elementos de $A$
pertenecientes al ideal maximal de $A$. Si $\mathfrak{K}$ es el ideal
engendrado por estos últimos y $\mathfrak{L}=\mathfrak{I}+\mathfrak{K}$, se
tiene $\mathcal{O}_{Y',x}=A/\mathfrak{L}$ y, como se está en la situación de
(16.9.13.5), el homomorfismo canónico
\[
  \mathfrak{I}/\mathfrak{I}\mathfrak{L}
  \longrightarrow(\mathfrak{I}+\mathfrak{L}^2)/\mathfrak{L}^2
\]
es biyectivo. Pero esto muestra que la sucesión
\[
  0\longrightarrow\mathfrak{I}/\mathfrak{I}\mathfrak{L}
  \longrightarrow\mathfrak{L}/\mathfrak{L}^2
  \longrightarrow(\mathfrak{L}/\mathfrak{I})/
    (\mathfrak{L}/\mathfrak{I})^2\longrightarrow0
\]
es exacta (véase la demostración de (16.2.7)), y los módulos que figuran en
esta sucesión son precisamente los tallos en $x$ de los haces de (16.9.13.1).
La segunda afirmación resulta de que $\mathcal{N}_{Y'/Y}$ es un
$\mathcal{O}_{Y'}$-módulo localmente \emph{libre} (16.9.8) y de Bourbaki,
\emph{Alg.}, cap.~II, 3\textsuperscript{a} ed., §~1, n\textsuperscript{o}~11,
prop.~21.
